% \appendix

% \section{Case Study: Build Environment Extraction and Matched-Toolchain Rebuild}
% \label{appendix:rpds-rebuild}

% We present a detailed case study of \texttt{rpds-py~0.30.0}, a compiled
% Rust extension for persistent data structures that appeared in the
% \emph{content\_diff} bucket of our criticality sample.
% This case illustrates both the feasibility and the remaining limits of
% artifact-driven build environment recovery.

% \subsection*{Information Sources in the Artifact}

% The released wheel
% \texttt{rpds\_py-0.30.0-cp313-cp313-manylinux\_2\_17\_x86\_64.whl}
% exposes four independent channels of build environment information
% without requiring access to any CI logs or build metadata outside the
% artifact itself.

% \begin{enumerate}
%   \item \textbf{Wheel filename.} The platform tag
%     \texttt{manylinux\_2\_17\_x86\_64} specifies the minimum glibc ABI
%     (2.17) and architecture (x86\_64), uniquely identifying the
%     \texttt{manylinux2014} container family as the build environment.

%   \item \textbf{\texttt{WHEEL} metadata.} The \texttt{Generator} field
%     records \texttt{maturin (1.10.2)}, pinning the exact build-tool
%     version.

%   \item \textbf{ELF \texttt{.comment} section.} Running
%     \texttt{readelf~-p~.comment} on the bundled shared object reveals
%     the full compiler chain embedded by the linker:
%     \begin{itemize}
%       \item \texttt{rustc 1.91.1 (ed61e7d7e 2025-11-07)}
%       \item \texttt{GCC: (GNU) 10.2.1 20210130 (Red Hat 10.2.1-11)}
%       \item \texttt{Linker: LLD 21.1.2}
%     \end{itemize}
%     The GCC version (10.2.1) independently corroborates the
%     \texttt{manylinux2014} container, whose base image ships exactly
%     that compiler.

%   \item \textbf{Source repository.} Fetching \texttt{pyproject.toml} at
%     the attested commit (\texttt{c38c9795}, recovered from the Fulcio
%     X.509 certificate) confirms \texttt{build-backend = "maturin"} with
%     no \texttt{rust-toolchain.toml} present for this release.
% \end{enumerate}

% \subsection*{Matched-Toolchain Rebuild}

% Using the information above we performed a controlled rebuild inside a
% fresh \texttt{quay.io/pypa/manylinux2014\_x86\_64} container:
% \texttt{rustup install 1.91.1} pinned the Rust toolchain, and
% \texttt{pip install maturin==1.10.2} pinned the build tool.
% We then ran \texttt{maturin build --release} at commit \texttt{c38c9795}.

% The resulting \texttt{.comment} section was byte-for-byte identical to
% the original:

% \begin{quote}
% \small
% \texttt{GCC: (GNU) 10.2.1 20210130 (Red Hat 10.2.1-11)}\\
% \texttt{Linker: LLD 21.1.2 (/checkout/src/llvm-project/llvm 8c30b9c...)}\\
% \texttt{rustc version 1.91.1 (ed61e7d7e 2025-11-07)}
% \end{quote}

% Every section header, symbol count (\texttt{.symtab}: \texttt{0x12300}
% entries), and code section (\texttt{.text}: \texttt{0x086adf} bytes)
% matched between the two wheels.
% The overall archive size converged from a 281~KB discrepancy (when
% \texttt{-{}-strip} was incorrectly applied) to a 232-byte residual.

% \subsection*{Residual Nondeterminism: Cargo Registry Path Hashes}

% Despite identical toolchain, build tool, container, commit, and build
% flags, the SHA-256 of the rebuilt wheel differed from the original.
% Section-level comparison isolated the divergence to three small regions:

% \begin{center}
% \begin{tabular}{lrrl}
% \toprule
% Section & Original & Rebuilt & Delta \\
% \midrule
% \texttt{.rodata}       & 32,416 B & 32,360 B & $-$56 B \\
% \texttt{.eh\_frame}    & 64,308 B & 64,312 B & $+$4 B  \\
% \texttt{.relro\_padding} & 696 B  &  744 B   & $+$48 B \\
% \texttt{.strtab}       & 206,566 B & 206,352 B & $-$182 B \\
% \bottomrule
% \end{tabular}
% \end{center}

% The \texttt{.strtab} divergence is caused by Cargo embedding absolute
% paths to cached \texttt{.rlib} files in the symbol table.
% These paths include the Cargo registry index hash, e.g.,
% \texttt{/root/.cargo/registry/src/index.crates.io-\textit{HASH}/rpds-1.2.0/}.
% The registry index hash is derived from the registry URL and the
% snapshot timestamp of the local \texttt{crates.io} index, which differs
% between the original CI run and our fresh clone.
% This is a known source of build nondeterminism in Cargo that is
% independent of the source code, toolchain, or build flags~\cite{cargo-nondeterminism}.

% \subsection*{Implications}

% This experiment demonstrates that the artifact itself provides sufficient
% information to reconstruct the build environment for compiled Rust
% extensions (\textbf{RQ4}): toolchain, container, and build tool can all
% be recovered without CI logs.
% The remaining barrier is Cargo's registry-path embedding, which would
% require either (a) a \texttt{CARGO\_REGISTRY\_TOKEN}-equivalent cache
% snapshot from the original build, or (b) upstream support for
% \texttt{CARGO\_HOME} normalisation analogous to
% \texttt{SOURCE\_DATE\_EPOCH} for timestamps.
% Until such normalisation is adopted, compiled Rust wheels face an
% irreducible nondeterminism floor that prevents byte-exact reproduction
% even under matched toolchains.

% % ─────────────────────────────────────────────────────────────────────────────

% \section{Case Study: Container Environment Mismatch in C-Extension Wheels}
% \label{appendix:pillow-rebuild}

% While the \texttt{rpds-py} case study isolates Cargo-level nondeterminism
% under a matched toolchain, the dominant failure mode for compiled wheels in
% our criticality sample is structurally earlier: the verifier's rebuild runs
% on the native host rather than inside the manylinux container the publisher
% used.
% We illustrate this with \texttt{Pillow~12.1.1}, a widely used C imaging
% library (\emph{cf.} the \texttt{preshed}, \texttt{srsly}, and
% \texttt{scikit-learn} cases, which exhibit the same root cause).

% \subsection*{Observed Divergence}

% The published wheel carries the platform tag
% \texttt{manylinux2014\_x86\_64} / \texttt{manylinux\_2\_17\_x86\_64},
% indicating it was produced inside a \texttt{quay.io/pypa/manylinux2014\_x86\_64}
% container (glibc $\geq$ 2.17, GCC 10.2.1).
% Our pipeline rebuilt Pillow natively on a Fedora~42 host (GCC~15.2.1),
% producing a wheel tagged \texttt{linux\_x86\_64}.
% The platform tag divergence is visible without even inspecting the binaries.

% Extracting the ELF \texttt{.comment} section from
% \texttt{\_avif.cpython-313-x86\_64-linux-gnu.so} in the published wheel
% confirms:

% \begin{quote}
% \small\texttt{GCC: (GNU) 10.2.1 20210130 (Red Hat 10.2.1-11)}
% \end{quote}

% Our rebuilt \texttt{.so} files, compiled with GCC 15.2.1, differ in
% every section that contains compiler-generated code: \texttt{.text},
% \texttt{.rodata}, \texttt{.eh\_frame}, and \texttt{.data.rel.ro}.
% All 8~shared objects in the wheel differ; only pure-Python files and
% \texttt{METADATA} match.

% \subsection*{Compounding Factor: Bundled Libraries}

% manylinux wheels statically bundle system libraries (libjpeg, libpng,
% libwebp, libtiff, libfreetype) that are not present on a standard host.
% Pillow's CI uses \texttt{auditwheel repair} to embed patched copies of
% these libraries into the wheel before upload.
% A native build links against the host's dynamic libraries, producing a
% structurally different binary even if the Pillow source and GCC version
% were somehow matched.
% This means that for manylinux wheels with bundled dependencies,
% \emph{matching the container alone is insufficient}: the exact
% \texttt{auditwheel} invocation and library versions must also be reproduced.

% \subsection*{Cython Source Nondeterminism: \texttt{srsly~2.5.3}}

% A related but distinct failure appears in \texttt{srsly~2.5.3}
% (9~content-differing files), which uses Cython to compile msgpack and
% ujson bindings.
% Beyond the expected \texttt{.so} divergence, the published wheel ships
% pre-generated Cython C++ sources (\texttt{\_packer.cpp},
% \texttt{\_unpacker.cpp}, \texttt{\_epoch.cpp}) directly in the wheel
% archive.
% These files differ from our rebuild because the installed Cython version
% on the host generates different C++ output than the version used during
% the original release build.
% The wheel thus embeds two layers of nondeterminism: the compiled object
% code and the intermediate C++ source it was compiled from.

% \subsection*{Implications}

% Together, \texttt{Pillow}, \texttt{preshed}, \texttt{scikit-learn}, and
% \texttt{srsly} represent a failure class distinct from the Cargo hash
% nondeterminism in \texttt{rpds-py}.
% The cause here is not residual build-system nondeterminism under a matched
% environment but a \emph{missing container layer}: the verifier pipeline
% did not detect that the target wheel was built inside a manylinux container
% and therefore did not reproduce that container before rebuilding.
% The fix — inferring the required container from the wheel's platform tag
% and the published \texttt{.so}'s \texttt{.comment} section (as
% \texttt{extract\_build\_env.py} does) — is well-defined and automatable.
% Until container-aware rebuild is integrated into the pipeline, all
% manylinux-tagged compiled wheels will fall into the
% \emph{other\_content\_mismatch} bucket regardless of source availability.
